\section{Preparation, timing, and causal retention}
\label{sec:foundation-budget}

The purpose of this section is to specify the common operands for the next
source-only decoder and detector-to-record calculations. Numerical laboratory
values are not assigned to an unspecified preparation, control, or memory.
The estimates below apply to explicitly declared channels and states and can
be filled with source-derived bounds as those constructions are completed.

\subsection{A preparation specification with retained failure}
Let $\rho_*=|\Psi_0\rangle\langle\Psi_0|$ denote the exact Chapter VI
hard-eigenbranch/soft-vacuum preparation in one declared lane. A physical
preparation must return a complete instrument with a ready flag and all
failure flags. For each operational input label $j$, write
\begin{equation}
 \widehat\rho_j=(1-f_j)|{\rm ready}\rangle\langle{\rm ready}|\otimes\rho_j^{\rm r}
       +f_j|{\rm fail}\rangle\langle{\rm fail}|\otimes\rho_j^{\rm f}.
 \label{eq:foundation-ready-state}
\end{equation}
Multiple failure outcomes can be retained as separate blocks. Specify the
ready-branch trace-distance error
\begin{equation}
 \delta_{{\rm prep},j}=\tfrac12\|\rho_j^{\rm r}-\rho_{*,j}\|_1.
\end{equation}
When $\rho_{*,j}$ is pure, a fidelity floor
$\langle\Psi_{*,j}|\rho_j^{\rm r}|\Psi_{*,j}\rangle\ge1-\zeta_j$
gives $\delta_{{\rm prep},j}\le\sqrt{\zeta_j}$ by the standard
trace-distance/fidelity relation~\cite{Watrous}.
This requires fidelity with the full specified source--field preparation;
a source marginal fidelity alone does not certify the joint field state.

If failure and abstention are counted as unsuccessful decoding, the complete
attempt law obeys the exact identity
\begin{equation}
 e_j^{\rm all}=f_j+(1-f_j)e_j^{\rm r}.
 \label{eq:foundation-failure-error}
\end{equation}
Consequently $f_j\le\bar f$ and $e_j^{\rm r}\le\bar e$ imply
$e_j^{\rm all}\le\bar f+(1-\bar f)\bar e$.
Conditioning on the ready flag defines a different, explicitly heralded
protocol. Its success probability and retained failed attempts remain part
of the resource account. A single experiment must choose whether it reports
the all-attempt or ready-conditional record criterion.

\subsection{A channel error budget that permits correlations}
For each label $j$, let $\Phi_k$ and $\widetilde\Phi_k$ be reference and actual
trace-preserving channels for the successive writing, switching, storage,
transport, access, and decoding stages. Classical flags are included in their
outputs. Registers needed to retain environmental correlations are included
in the stage input. The allowed states must be rich enough that each bound
below holds for the states on which that stage is actually used; if an
arbitrary reference is retained, use the corresponding bound with that
reference. Define
\begin{equation}
 \delta_k=\sup_{j,t,\omega\in\mathcal S_{k,j,t}}
  \tfrac12\|[(\widetilde\Phi_k-\Phi_k)\otimes I](\omega)\|_1.
 \label{eq:foundation-stage-error}
\end{equation}
The time supremum is over the declared receiving interval. A half-diamond-norm
bound is sufficient; a bound uniform on the relevant reachable states also
suffices. An environment with memory is retained as part of the joint state;
independent noise between stages or cells is not assumed.

\begin{proposition}[Composed readout error]
If the reference decoder has worst-label error at most $e_*(t)$ and all
stage bounds apply on the declared time interval, then the actual ready-branch
decoder satisfies
\begin{equation}
 \sup_{j,t} e_j^{\rm r}(t)
 \le \min\{1,\ \sup_t e_*(t)+\delta_{\rm prep}+\sum_k\delta_k\}.
 \label{eq:foundation-error-budget}
\end{equation}
\end{proposition}
\begin{proof}
Insert the reference preparation and each reference channel one at a time
into the composition. The triangle inequality bounds the final trace
distance by $\delta_{\rm prep}+\sum_k\delta_k$; subsequent trace-preserving
channels contract each difference. The success effect has spectrum in
$[0,1]$, so its probability changes by at most that trace distance. Apply the
bound to every label and time, then take the supremum. This argument also
applies to a channel whose output contains both a classical instrument label
and the corresponding quantum state.
\end{proof}
Equation~\eqref{eq:foundation-failure-error} then adds retained preparation
failure when the all-attempt criterion is used. Do not count that same
failure a second time inside an all-attempt channel-distance term.

\subsection{Timing errors and the energy account}
All internal ages use the Chapter I dimensionless chart. Let the actual
readout ages differ from the nominal ages by $\xi_1,\xi_2$. If
$|\xi_j|\le a_j$, the overlap time difference changes by at most
$a_1+a_2$. The centered-amplitude estimate in the preceding overlap section
uses the full prepared energy variance. It is insensitive to a common
scalar energy reassignment only after the compensating external action
phase is included. Absolute fringe phase requires the matched hard mean
energy and the scalar control-action convention. A gap interval alone does
not fix that absolute phase.

For a switched or noncommuting control Hamiltonian, replacing the entire
experiment by an age shift is a separate assertion. Instead, a common-domain
strong Duhamel estimate can bound a declared perturbation $\Delta K_j(t)$:
\begin{equation}
 r_j\le\int_0^{t_f}
       \|\Delta K_j(u)\Psi_{{\rm ref},j}(u)\|\,du.
 \label{eq:foundation-control-forcing}
\end{equation}
The forcing must exist and be integrable in Hilbert norm. The reference
orbit and domain assumptions must be stated for the actual singular-source
control law. A nominal pulse duration by itself supplies no value of $r_j$.
For normalized pure outputs, the resulting state trace distance is at most
$\min\{1,r_j\}$, so it contributes directly to the decoder error budget.
For branch overlap it contributes $r_1+r_2$ by the triangle inequality.

If the good-timing event has probability at least $1-\beta$ for every label,
and the conditional good-event decoding error is at most $e_{\rm good}$,
then the unconditional error is bounded by
\begin{equation}
 e\le\beta+(1-\beta)e_{\rm good}.
 \label{eq:foundation-timing-tail}
\end{equation}
The conditional input on that event must satisfy the preparation and control
bounds used for $e_{\rm good}$. Independence between timing, preparation,
and detector noise is not needed when the bounds are uniform or formulated
on their joint conditioned state.

\subsection{The specification to carry into the next calculation}
The following ledger distinguishes a defined mathematical operand from an
unassigned implementation value.
\begin{center}\small
\begin{tabularx}{\linewidth}{@{}l >{\raggedright\arraybackslash}X >{\raggedright\arraybackslash}X@{}}
\toprule
Operand & Meaning & Present input or next receiving construction\\
\midrule
$\rho_*,\Omega,M,v_s^{\rm up}$ & Joint reference and certified bounds
 & Retained Chapter VI/Chapter I package; exact hard branches are defined.\\
$\bar E,\Phi_0,\Phi_1$ & Absolute energy and explicit hard state operands
 & Recover certified hard-state data for numerical absolute phase and preparation.\\
$f_j,\delta_{{\rm prep},j}$ & Failure probability and ready-state accuracy
 & Bound from a complete finite preparation instrument.\\
$a_j,\beta,r_j$ & Age uncertainty, timing tail, control mismatch
 & Derive from the specified clock, switch, and actuator.\\
$\mathcal N_t$ & Storage including correlated noise and read disturbance
 & Choose a physical memory and retained environment.\\
$\mathcal A_{\mathcal D,t}$ & Transport and permitted receiving access
 & State receiver region, causal arrival, and available controls.\\
$F_j(t),e_*(t)$ & Fixed experiment's decoder and reference error
 & Next source-only quadrature benchmark, then actual detector bins.\\
$\mathcal T_{\rm mem},\mathcal T_{\rm read}$ & Memory and receiving windows
 & Use a causal onset and uniform retention/readout bounds.\\
\bottomrule
\end{tabularx}
\end{center}

For repeated reads, the channel sequence includes every earlier read and its
retained output. Uniform single-read performance alone specifies no bound
on a sequence of disturbed reads. Initial source--apparatus correlations,
finite control references, and failure recipients enter the declared state
and channels rather than being discarded at the composition step.

The interface to the later two-arm transport task is fixed now: one full
source--field state, declared branch controls and receiving regions, both
polarizations, all source sectors, even energy and counterterms, and the
same preparation/clock chart. The next detector calculation can then return
a complete conditional law and its error budget under the stated transport
law; actual two-arm transport will supply its additional controlled
discrepancy.

